% !TeX root = ../../Book2.tex
\chapter*{后记}
\addcontentsline{toc}{chapter}{后记}
\markboth{后记}{后记}

本卷从模出发，几次改变了问题中“同一个”的含义。给定主理想整环上的有限生成模，我们可以换生成元、换关系，最后用不变因子辨认其同构类型；给定线性算子，允许的是同一个空间的换基，相应的矩阵按相似变换。双线性型的矩阵却按合同变化，二次型的 Witt 类还进一步忽略双曲平面。写在纸上的矩阵都可能改变，允许怎样改变、改变后留下什么，才是每次分类的起点。

\ProseParagraphBreak
不变量也有强弱之分。特征多项式相同，算子未必相似；合成因子及其重数相同，有限长度模也未必同构。第十二章的不可分解直和分解回答的是另一种问题，不能由一串合成因子直接拼出。一个计算得到的不变量若能区分两个对象，便已很有用；若要声称它完成分类，还须证明相同数据足以得到所指定的等价。模的分类要求构造模同构，型的分类要求构造保持配对的线性同构，模范畴的比较则要求给出函子及其复合与恒等函子之间的自然同构。

\ProseParagraphBreak
这些比较所用的映射，在前几章已有具体的产生方式。自由模中的基像、商模中被杀掉的关系，以及张量积中的平衡关系，都决定了哪些映射能够唯一延拓或下降。第五章的自然性进一步要求这种对应与原对象的映射相容；第六章的 Hom--张量伴随将一侧的线性映射改写成另一侧的线性映射，第七章再分别检查 Hom 与张量函子何时保持正合性。这些构造到了 Morita 理论中组合成互逆的张量与 Hom 函子，比较的对象也从单个模变为模及其全部同态。

\ProseParagraphBreak
Morita 等价讨论含幺环的模范畴是否等价；它保留简单模、正合列等范畴结构，却不要求两个环的底层加法群或向量空间维数相同。Brauer 等价则固定域 \(k\)，只比较中心单 \(k\) 代数，并要求有关同构尊重 \(k\) 的标量作用。矩阵大小可以变化，但不能因此把这两种等价不加条件地混用；对一般基环上的 Azumaya 代数，零类的代表也未必是矩阵环，而可以是忠实有限投射模的自同态代数。

\ProseParagraphBreak
四个附录把这些结论带到不同的计算与结构问题中。附录 A 用精确矩阵计算和可核验的证书检验标准形；附录 B 从二次型的 Clifford 关系构造代数，再研究对合及 Spin 群。附录 C 在底模忠实、有限生成且投射的假设下，用全部极大纤维检测 Azumaya 性，并区别底模的局部自由性、代数的分裂与整体 Brauer 类。附录 D 则在代数关系之外加入范数和算子拓扑：同一组生成元取范数闭包或弱闭包，可能得到不同的代数；群全与约化完备化之间的满射也可能有非零核。

\ProseParagraphBreak
这些理论也各有止步之处。第四章的 Smith 标准形依赖主理想整环等条件；Witt 类有意忽略所加的双曲平面，单凭这个类不能恢复原型的维数；Brauer 群记录中心单代数的稳定差别，却不替我们解出每个具体的范数方程。第六章已经构造模的局部化，第三卷还要比较各个素理想处的局部信息如何反映整体结构，并研究完备化。第四卷用消解和导出函子研究正合性失效及扩张的来源；第五卷从本卷已有的群与 Lie 代数作用继续研究表示的分解和相应的上同调。面对这些新问题，本卷所训练的判断仍然具体：先确定标量、生成元与关系，再说明允许怎样比较对象，以及所得数据究竟能恢复什么。

\bigskip
\begin{flushright}
R. EverStarine

2026年10月1日
\end{flushright}
